\section{Related Work}
\label{sec:related}

Our work sits at the intersection of four threads of prior research: (1) AI and machine
learning for infectious-disease modelling, (2) automated or partially automated construction
of epidemiological models, (3) uses of LLMs in epidemiology, and (4) model-driven
engineering support for epidemiological modelling. The key distinction of our framework is
that it does not treat the LLM as a forecasting model or as a free-form scientific assistant.
Instead, it treats the LLM as an \emph{extraction and reconstruction instrument} whose output is checked against the paper text, a curated reference model, and structural consistency rules.

\subsection{AI and Machine Learning for Infectious-Disease Modelling}

Recent reviews show that AI is becoming a broad methodological layer in infectious disease epidemiology rather than a single technique. Liu \emph{et al.}~\cite{liu2023machinelearninginfectiousdisease} survey machine-learning approaches for infectious disease risk prediction and organize the literature into statistical prediction, data-driven machine learning, and epidemiology-inspired machine learning. Kraemer \emph{et al.}~\cite{Kraemer2025AI_epidemics} further broaden the picture by positioning AI for infectious-disease epidemics as a combination of machine learning, computational statistics, information retrieval, and data science, with applications spanning surveillance, data integration, forecasting, and decision support. Chopra \emph{et al.}~\cite{chopra2023differentiableagentbasedepidemiology} introduce differentiable agent-based epidemiology, showing how mechanistic simulators can be made amenable to gradient-based learning and linked to auxiliary data sources.

\subsection{LLMs in Epidemiology}

Wibaek \emph{et al.}~\cite{wibaek2023llm_epidemiology}
show that large language models can be applied to epidemiological cohort text to build
predictive models for downstream health outcomes. Matthay \emph{et al.}~\cite{Matthay2025}
argue more broadly that AI can assist multiple stages of causal research in epidemiology,
including data collection, variable construction, and decision support. Plecko \emph{et al.}~\cite{plecko2025epidemiologyllm} propose a benchmark for
observational distribution knowledge across real-world domains, including health-related
settings. Although this work is not about compartmental modelling, it is directly relevant to
our problem formulation: it cautions against assuming that an LLM internally ``knows'' the
population distributions or domain regularities needed for faithful epidemiological reasoning.
That observation supports our design choice to rely on paper-grounded extraction and explicit
validation rather than unconstrained generation.

\subsection{LLMs for Epidemic and Transmission Modeling}

The literature most closely related to our paper studies LLMs as assistants for epidemic or
transmission modeling. Kwok \emph{et al.}~\cite{KWOK20243254} present a case study in which
ChatGPT collaborates with a public-health practitioner to iteratively generate, refine, and
debug a transmission model, ultimately fitting prevalence data and estimating epidemic
quantities such as $R_0$ and final epidemic size. Their study shows
that LLMs can lower the technical barrier to building a working infectious-disease model. EpiLLM~\cite{gong2025epillmunlockingpotentiallarge}
recasts epidemic forecasting as an LLM adaptation problem, using spatio-temporal prompt
learning, mobility signals, and autoregressive next-token prediction to improve multi-step
forecasting accuracy on COVID-19 data. EpidemIQs~\cite{samaei2026epidemiqsprompttopaperllmagents}
proposes a multi-agent prompt-to-paper framework that can conduct literature review,
analytical derivation, network and mechanistic modelling, stochastic simulation,
visualization, and manuscript drafting. This is highly relevant to our motivation because it
shows growing interest in end-to-end LLM-supported epidemic-model development.

%Our own earlier work~\cite{11344468} is the closest direct precursor. That study examined
%whether LLMs could generate plausible epidemiological models from textual descriptions in a
%single-case setting. The current paper goes beyond plausibility. We introduce a full
%paper-to-model pipeline, evaluate it across 10 diseases, compare LLM backends systematically,
%and, most importantly, define a reproducibility-oriented notion of \emph{completeness} that
%separates paper promises, extracted content, and reference agreement.

%\subsection{Model-Driven Engineering and Reproducible Model Representations}

%Model-driven engineering provides the formal substrate that makes our evaluation possible.
%MDE uses explicit metamodels to define the structure of domain artifacts and to support
%artifact generation, checking, and transformation~\cite{brambilla2012model}. In the
%infectious-disease setting, EpiMDE~\cite{epimde2024} is the most relevant prior platform: it
%proposes an extensible metamodel and integrated environment for epidemiological
%compartmental models, enabling structured model creation, simulation, and cross-model
%comparison. Our work directly builds on this idea of a formal epidemiological representation.
%The difference is that EpiMDE primarily supports \emph{expert-authored} models, while our
%framework targets \emph{paper-derived} models and studies whether LLM-assisted extraction can
%populate such a representation completely and correctly.

%\subsection{Retrieval, Matching, and Verification}

%RAG methods ground LLM outputs in retrieved evidence rather than relying only on
%parametric memory~\cite{lewis2020retrieval}. This is especially important in our setting,
%where many target items are short technical expressions, parameter symbols, or local textual
%commitments that appear only within the paper being processed. For that reason, our system
%uses paper-grounded retrieval over text chunks and parameter tables, not open-domain
%retrieval.

%Likewise, evaluating extracted epidemiological entities requires more than exact string match.
%Compartment and parameter names vary across papers through abbreviations, Greek/Latin
%notation changes, age qualifiers, and disease-specific terminology. We therefore adopt a
%domain-aware fuzzy matching strategy based on normalized labels and
%\texttt{SequenceMatcher}~\cite{python_difflib}. This evaluation layer is not merely an
%implementation detail; it is part of the contribution because it enables a fair measurement of
%whether two model artifacts refer to the same epidemiological concept.

%Overall, the literature now contains substantial work on AI for epidemiology, LLM-supported
%epidemic modelling, forecasting, and model automation. What remains underexplored is the
%specific reproducibility question addressed here: given only a published modelling paper, how
%faithfully can one recover the formal compartmental model it actually specifies, and how
%should that faithfulness be measured? Our framework fills that gap.